\section{Two accounts of what demotion costs}
\label{sec:theory}

This section fixes the estimand (\S\ref{sec:estimand}), then states two
accounts of it in those same units: a training-gap interpretation built
on the spectral account of the inference-side literature
(\S\ref{sec:spectral}), and ours
(\S\ref{sec:ouraccount}). They are not variants
of one model: they disagree about which quantity governs a route's
cost, about what happens at pure noise, and about whether a single
parameter can order all routes.

\subsection{The estimand}
\label{sec:estimand}

\paragraph{Objective and noising.}
A flow-matching model regresses the constant-velocity transport between
a clean latent and pure noise. With clean latent $x$, noise
$\epsilon \sim \mathcal{N}(0, I)$, and noise level $\sigma \in (0,1]$,
the noised input is $z_\sigma = (1-\sigma)\,x + \sigma\,\epsilon$ and
the target is the velocity $v = \epsilon - x$; with caption condition
$c$ the loss is $\|\hat v_\theta(z_\sigma, \sigma, c) - v\|^2$.

\paragraph{Demotion.}
We write $e_0 {\to} e$ for a \emph{demotion}: the native latent at the
route's source edge $e_0$ is replaced with a coarser-grid one at its
target edge $e$.
With $y$ the image at its native resolution, $\mathcal{E}$ the VAE
encoder, and $\mathcal{R}_e$ the pixel-space downscale to the tier-$e$
bucket (native aspect ratio preserved),
$x_{\src} = \mathcal{E}(y) \longrightarrow
x_{\dem} = \mathcal{E}\big(\mathcal{R}_{e}(y)\big)$ --- a latent with
proportionally fewer spatial tokens, the downscale acting in pixel
space, not latent space. The demoted step trains on
$z_\sigma = (1-\sigma)\,x_{\dem} + \sigma\,\epsilon$ with target
$v = \epsilon - x_{\dem}$, so the substitution changes both what the
network sees and what it is asked to predict, but only through the
spatial grid.

\paragraph{The estimand.}
Three \emph{arms} recur throughout: the source arm $\src$, which
trains on the native-resolution latent $x_{\src}$; the demoted arm
$\dem$, which trains on $x_{\dem}$; and the control arm $\reenc$,
which decodes and re-encodes the native latent at native resolution.
For an arm
$a \in \{\src, \reenc, \dem\}$, let
$\bar g_a(\sigma) = \mathbb{E}_\epsilon[\nabla_\theta \mathcal{L}]$ be
the population per-bin mean adapter gradient of a query (image,
caption). The \emph{population demotion distance} of the route,
$d_{e_0 \to e}$, is
\begin{equation}
d_{e_0 \to e}(\sigma) \;=\; 1 - \cos\!\big(\bar g_{\src}(\sigma),\,
\bar g_{\dem}(\sigma)\big)
\;=\; \tfrac{1}{2}\,\big\|\hat g_{\src} - \hat g_{\dem}\big\|^2
\;\ge\; 0,
\label{eq:gapdef}
\end{equation}
with $\hat g$ the unit-normalized gradients. Route-level scalars carry
the \emph{route} as a subscript, abbreviated to the target edge alone
($d_e$) wherever the source is fixed by context; the
control's distance $d_{\reenc}$ carries its arm label instead, since it
changes no grid (full convention: Appendix~\ref{app:notation}). The
\emph{reported} quantity of every bin-resolved map and every fit in this
paper is the \emph{re-encoding excess}
\begin{equation}
\gap_e(\sigma) \;:=\; d_e(\sigma) \;-\; d_{\reenc}(\sigma),
\label{eq:excess}
\end{equation}
the demotion distance beyond what the control already pays; $\gap_e$
can be negative, $d_e$ cannot. Endpoint and floor tables report
debiased estimates of $d_e$ with the control as its own row. Every
map in the body aggregates per example (cosine per image, then mean);
the batch-aggregate object SGD follows in batched training is a
distinct estimand, treated in Appendix~\ref{app:notation} --- where
measured, it reads as the \emph{more forgiving} of the two
(\S\ref{sec:map}).

\subsection{The spectral account}
\label{sec:spectral}

The spectral argument of \S\ref{sec:related} supplies the field's
working model of the same gap, and it is worth deriving in the units
Eq.~\ref{eq:gapdef} just fixed. Its strongest available form is the
tolerance-parameterized crossover of \citet{xiao2026spd}.

\paragraph{Demotion deletes a band.}
Write $u^{(\omega)}$ for the frequency-$\omega$ component of a field
$u$ in the latent's spatial Fourier decomposition, with $\omega$ in
normalized units --- cycles per sample of the \emph{native} grid, so
the native Nyquist sits at $\tfrac12$. The downscale $\mathcal{R}_e$
resamples the image with $e/e_0$ times fewer samples per axis, and a
grid at that density can represent no frequency above half its own
sampling rate; the resampler's anti-aliasing prefilter removes, rather
than aliases, whatever lies above. Idealizing the VAE round trip away
--- its residue is exactly what the $\reenc$ control of
\S\ref{sec:estimand} pays --- demotion is an ideal low-pass:
\begin{equation}
x_{\dem}^{(\omega)} \;=\;
\begin{cases}
x^{(\omega)}, & |\omega| < \omega_e,\\[3pt]
0, & |\omega| \ge \omega_e,
\end{cases}
\qquad
\omega_e \;=\; \frac{1}{2}\,\frac{e}{e_0}\,,
\label{eq:lowpass}
\end{equation}
deleting exactly the bands above the coarse grid's Nyquist frequency
$\omega_e$ and leaving the rest intact.

\paragraph{Carried to the training gradient.}
The noising is linear, so it acts band by band, and under the premise
that each clean-latent band is independently Gaussian,
$x^{(\omega)} \sim \mathcal{N}(0, P_\omega)$ with $P_\omega$ the
band's clean-latent power, \citet{xiao2026spd} attach to each band an
activation time: a crossover $t_\omega$ past which the Bayes-optimal
per-band velocity prediction sits within a tolerance $\delta$ --- the
account's one free parameter --- of the pure-noise answer
$\epsilon^{(\omega)}$; the choice $\delta = (1+P_\omega)/2$ recovers
the often-quoted equal-power (input-SNR) crossover
$\sigma_{\mathrm{eq}} = \sqrt{P_\omega}/(1+\sqrt{P_\omega})$ as one
slice of the family. How $\delta$ is set in practice --- including a
zero-free-parameter anchor measured from our own pipeline --- is
treated in Appendix~\ref{app:spectral}; no conclusion below hinges on
the choice.
The destroyed bands of Eq.~\ref{eq:lowpass} all sit at or above
$\omega_e$; on a decaying spectrum the lowest of them is the most
powerful, so it crosses last, and the safe-substitution boundary for
the whole route sits at $t_{\omega_e}$. Under the premise that
demotion costs only what the destroyed bands still contribute to the
prediction, the cost has one place to land: whatever those bands
contribute to the residual, and through the backward pass to the
adapter gradient, perturbs $\bar g_{\dem}$ away from $\bar g_{\src}$.
The implied gradient gap is that contribution, gated at the Nyquist
band's activation time:
\begin{equation}
\gap_e^{\mathrm{spec}}(\sigma) \;=\;
S_e(\sigma)\,\mathbf{1}\!\big[\sigma < t_{\omega_e}\big],
\qquad
t_\omega \;=\;
\frac{1}{1 + \sqrt{\delta \,/\, \big(P_\omega\,(1 + P_\omega -
\delta)\big)}},
\label{eq:spd}
\end{equation}
where $S_e(\sigma)$ is the destroyed-band contribution to the distance
of Eq.~\ref{eq:gapdef}. $S_e$ can be computed. At the residual level
the premise is complete: the Gaussian posterior mean is the classical
per-band Wiener shrinkage \citep{wiener1949}, so the cross-grid
mean-residual mismatch is confined to the destroyed band, ordered
across routes by destroyed-band energy, and follows in closed form
from the power spectrum (Appendix~\ref{app:posterior}); carrying it
into the gradient units of Eq.~\ref{eq:gapdef} costs one calibrated
gain, since the map passes through the network's Jacobian, about which
a spectral model of the \emph{data} says nothing.
Appendix~\ref{app:spectral} instantiates exactly this pipeline when
the account is scored.

Our claim is therefore not that $S_e$ cannot be estimated --- it can,
one way or another --- but that a gap estimate built from it is
structurally insufficient for $\gap_e$: three limits are fixed by the
\emph{form} of Eq.~\ref{eq:spd}, before any gradient is measured and
however $S_e$ is computed, and each is a place the account can be
falsified without arguing about $\delta$.
\textbf{(a)~One parameter orders every route}: all boundaries move
together under the single tolerance, so the family is totally ordered.
\textbf{(b)~Ratio only}: the cut $\omega_e = \tfrac{1}{2}\,e/e_0$ of
Eq.~\ref{eq:lowpass} is a function of the downscale \emph{ratio} alone,
so the estimate is structurally blind to any absolute-size effect. \textbf{(c)~No floor}: for every route and
every $\delta$ the predicted gap vanishes above a finite boundary, so
no $\sigma$-independent term is expressible.

\subsection{Our account: what demotion perturbs, and its angular
reduction}
\label{sec:ouraccount}
\label{sec:reduction}

We start one level lower, from what the substitution touches
in the gradient itself. Differentiating the loss of
\S\ref{sec:estimand} gives, per draw and up to a constant,
$g = \nabla_\theta \tfrac12\|\hat v_\theta - v\|^2 = J^{\!\top} r$,
with $r = \hat v_\theta - v$ the prediction residual and
$J = \partial \hat v_\theta / \partial \theta$ the Jacobian through
which the compute graph turns a residual into an adapter gradient.
Demotion changes exactly the two factors of this product, and nothing
else about the map from (image, noise, caption) to adapter gradient:
the residual, because it changes the data the objective consumes, and
the Jacobian, because it changes the compute graph the gradient flows
through. The estimand is built from noise means, so the relevant
residual object is the \emph{mean residual} at fixed query,
$\bar r_a(\sigma) = \mathbb{E}_\epsilon[\hat v_\theta - v]$ --- the
same averaging that defines $\bar g_a$ in \S\ref{sec:estimand} --- a
velocity-shaped field with a definite $\sigma$-curve for each arm;
write $\Delta \bar r(\sigma)$ for its cross-arm mismatch
$\bar r_{\dem} - \bar r_{\src}$, and $\Delta J$ likewise for the
noise-mean Jacobians. Averaging each factor over draws and expanding
the product (Appendix~\ref{app:branches}), the perturbation
$\delta g = \bar g_{\dem} - \bar g_{\src}$ decomposes as
\begin{equation}
\delta g \;=\;
\underbrace{J^{\!\top}\Delta \bar r(\sigma)}_{\text{data branch }B_e}
\;+\; \underbrace{\Delta J^{\!\top}\, \bar r}_{\text{graph branch }C_e}
\;+\; \underbrace{\Delta J^{\!\top}\Delta \bar r + \cdots}_{\text{remainder }R_e}.
\label{eq:branches}
\end{equation}
$J_{\src}$ and $J_{\dem}$ act on different grids, so the branches are
not literal matrix products across arms; they are defined
\emph{operationally}, as the two independent perturbation directions in
the shared adapter-parameter space where both gradients live, and
\S\ref{sec:ourscored} isolates each by its own intervention.

Read against Eq.~\ref{eq:branches}, the spectral account of
\S\ref{sec:spectral} keeps the input-mediated part of the data
branch, gated at Nyquist, and sets everything else to zero --- no
target-mediated share, no graph branch ($\Delta J \equiv 0$), no
interaction --- and the three limits fixed there are these discarded
terms read back. The graph branch does not reference $\sigma$ at all,
since the grid change does not: token count, rotary phase density,
and sequence-length-dependent normalization change identically at
every noise level, so nothing about noise masking can remove them ---
limit~(c)'s inexpressible $\sigma$-independent term --- and, being a
property of the two grids rather than of the downscale map, it need
not be a function of their ratio (limit~(b)); a gap governed by two
terms rather than one admits no single-parameter ordering of the
family (limit~(a)). Within the data branch, only the
\emph{input}-mediated part shrinks with $\sigma$: perturbing the
image by $\delta x$ perturbs the residual by
$D_z f\,[(1{-}\sigma)\,\delta x] + \delta x$, and the second,
target-mediated term carries unit weight at every $\sigma$, so at
$\sigma{=}1$ --- where the input is pure noise on every grid and the
input-mediated part vanishes in distribution --- the endpoint is
still not data-free by construction: at the Bayes level the surviving
mean residual is exactly the image's deviation from the
caption-conditional prior mean, $x - \mathbb{E}[x \mid c]$
(Appendix~\ref{app:posterior}). Whether that survivor is
\emph{resolvable in the angular estimand} is an empirical question,
discharged in \S\ref{sec:ourscored}.

\paragraph{The angular reduction.}
It remains to read Eq.~\ref{eq:branches} in the units of the estimand.
The geometry is pure bookkeeping, so the displayed forms live in
Appendix~\ref{app:geometry}; what the argument needs is their content.
Writing
$u^{\perp} = u - (\hat g_{\src}^{\!\top}u)\,\hat g_{\src}$ for the
component orthogonal to $\bar g_{\src}(\sigma)$, split the perturbation
into a rescaling and a rotation:
$\kappa_\parallel = \hat g_{\src}^{\!\top}\delta g / \|\bar g_{\src}\|$
and $\kappa_\perp = \|\delta g^{\perp}\| / \|\bar g_{\src}\|$.
Substituting $\bar g_{\dem} = \bar g_{\src} + \delta g$ into
Eq.~\ref{eq:gapdef} gives, \emph{exactly},
$d_e = 1 - 1/\sqrt{1 + \kappa_{\mathrm{eff}}^2}$ with
$\kappa_{\mathrm{eff}} = \kappa_\perp/(1 + \kappa_\parallel)$
(Eq.~\ref{eq:exact}): the gap is a saturating function of a single
rotation-to-scale ratio --- only the orthogonal component rotates the
gradient direction, and the cosine estimand is \emph{blind} to a
parallel rescaling, a fact the endpoint probes of \S\ref{sec:ourscored}
lean on. Taylor-expanding at small $\kappa$ and substituting
Eq.~\ref{eq:branches} into the square gives the \emph{four-term angular
expansion} (Eq.~\ref{eq:fourterm}), derived as unconditionally as the
exact form itself: a data share
$S_e = \|B_e^{\perp}\|^2 / (2\|\bar g_{\src}\|^2)$, a graph share
$F_e = \|C_e^{\perp}\|^2 / (2\|\bar g_{\src}\|^2)$, a projected
interaction
$I_e = \langle B_e^{\perp}, C_e^{\perp}\rangle / \|\bar g_{\src}\|^2$,
and a remainder $R'_e$ collecting the remainder branch, the
beyond-quadratic terms, and the parallel-component correction. $S_e$
and $F_e$ are non-negative; $I_e$ can carry either sign, a fact that
matters in Appendix~\ref{sec:headtohead}. At the harshest measured
operating point ($d = 0.30$ at $512$, $\sigma{=}1$) the implied
perturbation is $\kappa_{\mathrm{eff}} \approx 1.02$ --- not small ---
so quadratic-order reads on that route are indicative and the exact
form carries the quantitative claims. Note that Eq.~\ref{eq:spd} is
Eq.~\ref{eq:fourterm} with $F_e = I_e = 0$ and $S_e$ gated at Nyquist.

\paragraph{From the expansion to a measurable form.}
The reported object is the excess over the control, whose own expansion
is degenerate --- the $\reenc$ arm changes no grid, so
$\Delta J_{\reenc} = 0$ and both terms built from it vanish, leaving
$d_{\reenc} \approx S_{\reenc} + R'_{\reenc}$. Subtracting term by term
therefore cancels only the re-encode component of the data share and
passes the graph share and interaction through untouched. Four named
assumptions, each tested by a designated probe in
\S\ref{sec:evidence}, then trade the exact bookkeeping (in which no
term is separately measurable per bin) for a form built from measurable
curves and per-route constants:
(i)~\textbf{small remainder}, $R'_e$ negligible over the claimed
domain;
(ii)~\textbf{negligible projected interaction},
$|I_e| \ll S_e + F_e$;
(iii)~\textbf{graph-relative stationarity},
$\|C_e^{\perp}\| / \|\bar g_{\src}\|$ $\sigma$-constant (numerator and
denominator share the residual's scale), so $F_e$ is flat in $\sigma$
--- write $\Floor_e$ for its value;
(iv)~\textbf{data-branch factorization},
$\|B_e^{\perp}\| \approx a_e\,\|\Delta \bar r(\sigma)\|$: a
$\sigma$-independent route amplitude times the \emph{norm} of the
cross-grid mean-prediction-residual mismatch of Eq.~\ref{eq:branches},
asserted to be route-independent and directly measurable as such
(\S\ref{sec:ourscored}); the measured curve is an \emph{empirical}
closure of the mean residual's exact posterior form
(Appendix~\ref{app:posterior}), not an unconstrained fitted shape.
The first two are not unstructured hopes: both discarded terms obey
exact Cauchy--Schwarz bounds in the retained shares
($|I_e| \le 2\sqrt{S_eF_e} \le S_e + F_e$, the remainder entering at
half power in its own share), so (ii) carries independent content only
near share crossover (Appendix~\ref{app:geometry}). Applied,
\begin{equation}
\gap_e(\sigma)
\;\approx\;
\underbrace{\tfrac{1}{2}\,a_e^2 \cdot \big(\|\Delta \bar r(\sigma)\| \,/\, \|\bar g_{\src}(\sigma)\|\big)^{2}}_{\text{data term}}
\;+\; \underbrace{\Floor_e}_{\text{graph term}},
\label{eq:twoterm}
\end{equation}
the \emph{two-term account}, stated on the reported excess. The
quadratic power is fixed by the cosine geometry \emph{locally}; far
from the small-mismatch regime the licensed read is the un-expanded
parent form, Eq.~\ref{eq:exact} with loading
$\kappa_e(\sigma) = a_e\, \|\Delta \bar r(\sigma)\|/\|\bar g_{\src}(\sigma)\|$ (the
same amplitude, recovering Eq.~\ref{eq:twoterm}'s data term in the
small-$\kappa$ limit), whose data term
\emph{saturates} instead of overshooting. The loading itself is an
\emph{empirical} ingredient, not a derivation, so
Appendix~\ref{sec:headtohead} scores both forms held out, labelling
Eq.~\ref{eq:exact} the best empirical predictor and Eq.~\ref{eq:twoterm}
the derived small-perturbation form. To keep the ledger explicit:
Eqs.~\ref{eq:exact} and~\ref{eq:fourterm} are unconditional;
Eq.~\ref{eq:twoterm} is conditional on the four assumptions; the
coefficients $a_e$, $\|\Delta \bar r(\sigma)\|$,
$\|\bar g_{\src}(\sigma)\|$ and $\Floor_e$ are measured. The account was pre-registered with an
outcome matrix --- including a kill switch, all endpoint gaps
${\approx}\,0$, that would have falsified the graph term --- before the
discriminating runs.

\paragraph{The graph term has identifiable parts.}
Changing the grid changes (i)~the rotary coordinate system, in which
every band of the position embedding accumulates phase at a different
density across the image, and (ii)~non-positional graph statistics:
attention softmax over a different token count,
sequence-length-dependent normalization, and the capacity of the coarse
graph to approximate the fine graph's computation (the conceptual frame
is attention as a discretization of a function-space operator,
\citealp{calvello2024continuum}, whose statistics need not be
discretization-invariant). Splitting
$\Floor_e = \mathrm{RoPE}_e + \mathrm{Resid}_e$ yields one sharp
prediction: the positional share is a \emph{coordinate-system
artifact}, so an exact positional-interpolation intervention should
remove $\mathrm{RoPE}_e$ wherever the data branch is inactive
($\sigma \to 1$); a band-counting sketch attaches an
absolute-\emph{size} governor to $\mathrm{Resid}_e$. Like YaRN's own
band thresholds \citep{peng2024yarn}, that sketch's constants are
calibrated, not derived; we claim no closed form for $\Floor_e$.
